\documentclass[12pt,a4paper]{article}

% ---------------------------------------------------------
% Packages
% ---------------------------------------------------------
\usepackage[a4paper,margin=1in]{geometry}
\usepackage[T1]{fontenc}
\usepackage{newtxtext,newtxmath}
\usepackage{setspace}
\usepackage{enumitem}
\usepackage{booktabs}
\usepackage{array}
\usepackage{graphicx}
\usepackage{float}
\usepackage{titlesec}
\usepackage{fancyhdr}
\usepackage[hidelinks]{hyperref}

% ---------------------------------------------------------
% Document formatting (following Journal 02)
% ---------------------------------------------------------
\setstretch{1.12}
\setlength{\parindent}{0pt}
\setlength{\parskip}{6pt}

\titleformat{\section}
{\large\bfseries}
{\thesection.}{0.5em}{}

\titleformat{\subsection}
{\normalsize\bfseries}
{\thesubsection}{0.5em}{}

\pagestyle{fancy}
\fancyhf{}
\fancyfoot[C]{\thepage}
\renewcommand{\headrulewidth}{0pt}

\begin{document}

% ---------------------------------------------------------
% Title Page
% ---------------------------------------------------------
\begin{titlepage}
\centering
\vspace*{0.8in}

{\Large\bfseries Smart City Civic Complaint and Management System\par}

\vspace{0.25in}

{\large\bfseries Software Engineering Project Journal -- 04\par}

\vspace{0.8in}

\textbf{Submitted to:}\par
Nisha Thakur\par
Thapar Institute of Engineering and Technology\par

\vspace{0.45in}

\textbf{Submitted by:}\par
Hittin Gupta [1024170437]\par

\vfill

\textbf{Journal 04}\par
Data Flow Diagram Design\par
November 2026

\end{titlepage}

% ---------------------------------------------------------
% 1. Objective
% ---------------------------------------------------------
\section{Objective}

The objective of this journal is to document the preparation of Level 0 and
Level 1 Data Flow Diagrams (DFDs) for the Smart City Civic Complaint and
Management System.

The main objectives were:

\begin{itemize}[leftmargin=*]
    \item To understand how information enters, moves through, and is stored
    within the complaint management system.
    \item To identify the external entities that interact with the system.
    \item To prepare a Level 0 DFD representing the system as a single
    high-level process.
    \item To prepare a Level 1 DFD by decomposing the system into its major
    internal processes.
    \item To represent the data exchanged between citizens, department
    officers, field workers, system processes, and the database.
    \item To improve the clarity of the system design before final
    documentation and evaluation.
\end{itemize}

% ---------------------------------------------------------
% 2. Work Completed
% ---------------------------------------------------------
\section{Work Completed}

During this phase, the existing system workflow and its major data
interactions were studied. Based on the complaint lifecycle, two Data Flow
Diagrams were prepared: a Level 0 Context Diagram and a Level 1 Process
Decomposition Diagram.

The diagrams were created to represent the movement of information between
the users, the system processes, and the PostgreSQL database.

The major activities completed were:

\begin{enumerate}[leftmargin=*]
    \item Identification of external entities interacting with the system.
    \item Identification of the main inputs and outputs exchanged with each
    external entity.
    \item Preparation of the Level 0 DFD.
    \item Decomposition of the system into five major processes for the
    Level 1 DFD.
    \item Representation of the PostgreSQL database as the central data store.
    \item Review of the data flows and workflow hand-offs represented in
    both diagrams.
    \item Addition of the completed diagrams to the project documentation.
\end{enumerate}

% ---------------------------------------------------------
% 3. Level 0 DFD
% ---------------------------------------------------------
\section{Level 0 Data Flow Diagram}

The Level 0 DFD, also known as the Context Diagram, represents the complete
Smart City Civic Complaint and Management System as a single process. It
shows the external entities that communicate with the system and the major
information exchanged between them.

The three external entities represented are:

\begin{itemize}[leftmargin=*]
    \item \textbf{Citizen:} Provides registration and login details,
    complaint information, category, location, photographs, and resolution
    feedback. The system returns authentication results, complaint
    acknowledgement, complaint ID, status, and resolution details.
    \item \textbf{Department Officer:} Provides login details, complaint
    verification or rejection decisions, remarks, and worker assignment
    requests. The system provides complaint lists, complaint details,
    history, available workers, and assignment results.
    \item \textbf{Field Worker:} Provides login details, task acceptance,
    work-start information, progress updates, work notes, and completion
    evidence. The system provides assigned task details and complaint
    information.
\end{itemize}


% ---------------------------------------------------------
% 4. Level 1 DFD
% ---------------------------------------------------------
\section{Level 1 Data Flow Diagram}

The Level 1 DFD expands the main system process into five major processes.
It provides a more detailed representation of how complaint-related data
is processed and stored.

The five processes represented are:

\begin{enumerate}[leftmargin=*]
    \item \textbf{1.0 User Authentication and Role Validation:}
    Handles user registration and login information and validates the
    user's role before providing access.
    \item \textbf{2.0 Complaint Submission and Evidence:}
    Receives complaint details, category, location, and supporting
    photographs from the citizen and records the submitted complaint.
    \item \textbf{3.0 Officer Review and Assignment:}
    Allows the department officer to review a complaint, provide a
    verification decision, and assign a field worker.
    \item \textbf{4.0 Worker Task Execution:}
    Handles task acceptance, work progress, work notes, and completion
    evidence submitted by the field worker.
    \item \textbf{5.0 Tracking, Resolution and Feedback:}
    Provides complaint status and history, supports citizen verification
    of completed work, and records resolution feedback.
\end{enumerate}

The Level 1 diagram also represents the PostgreSQL database as the central
data store. It contains information related to users, complaints,
assignments, status history, complaint evidence, and feedback. The
processes read from and write to this data store as required by their
operations.



% ---------------------------------------------------------
% 5. Outcome
% ---------------------------------------------------------
\section{Outcome}

The Level 0 and Level 1 Data Flow Diagrams were prepared to document the
system's data movement at two levels of abstraction. The Level 0 diagram
provides an overall view of the system and its external interactions,
whereas the Level 1 diagram explains the internal processes and their
interaction with the PostgreSQL data store.

These diagrams improve the understanding of the system's data flow and
provide supporting documentation for the overall architecture and
implementation of the Smart City Civic Complaint and Management System.

% ---------------------------------------------------------
% 6. Next Planned Work
% ---------------------------------------------------------
\section{Next Planned Work}

The next stage will focus on reviewing the completed project documentation,
checking the consistency of the diagrams with the implemented system, and
preparing the final project submission and demonstration.

\end{document}